\appendix

\chapter{Gitコマンド早見表}

\section*{この付録の使い方}
\addcontentsline{toc}{section}{この付録の使い方}

本書で解説してきたGitコマンドを、カテゴリ別に一覧で掲載しています。「あのコマンド、何だったかな」と思ったときに素早く引ける辞書としてお使いください。各コマンドの詳しい解説は、対応する章を参照してください。

また、末尾には「やりたいこと」から逆引きできる表も用意しています。

\section{初期設定（→ 第2章）}

Gitを使い始めるときに最初に行う設定です。\cmd{--global} は全リポジトリ共通の設定、省略するとそのリポジトリだけの設定になります。

\begin{lstlisting}[language=bash]
# ユーザー名の設定（コミットに記録される名前）
git config --global user.name "あなたの名前"

# メールアドレスの設定（コミットに記録されるメールアドレス）
git config --global user.email "your-email@example.com"

# デフォルトブランチ名をmainに統一する
git config --global init.defaultBranch main

# コミットメッセージ編集時に使うエディタを指定する
git config --global core.editor "code --wait"

# 現在の設定を一覧表示する
git config --list

# 特定の設定項目の値を確認する
git config user.name

# 設定ファイルをエディタで直接開く
git config --global --edit
\end{lstlisting}

\section{リポジトリの作成と取得（→ 第5章）}

\begin{lstlisting}[language=bash]
# 現在のディレクトリをGitリポジトリとして初期化する
git init

# リモートリポジトリの完全なコピーをローカルに作成する
git clone URL

# 最新のコミットだけを取得する（高速だが履歴が不完全）
git clone --depth 1 URL

# 特定のブランチを指定してクローンする
git clone -b ブランチ名 URL

# サブモジュールも含めてクローンする
git clone --recursive URL
\end{lstlisting}

\section{基本操作（→ 第6章）}

日常的に最も多く使うコマンド群です。add → commit → push の流れが基本サイクルになります。

\begin{lstlisting}[language=bash]
# 現在のリポジトリの状態を確認する（最も頻繁に使うコマンド）
git status

# 短縮形式で状態を確認する
git status -s

# 特定のファイルをステージングエリアに追加する
git add ファイル名

# 複数のファイルをまとめてステージングする
git add file1.py file2.py

# カレントディレクトリ以下の全変更をステージングする
git add .

# 削除されたファイルも含めて全変更をステージングする
git add -A

# 変更を部分的にステージングする（対話形式で選択）
git add -p

# ステージングした変更をコミットする（メッセージ付き）
git commit -m "コミットメッセージ"

# 複数行のコミットメッセージを書く
git commit -m "1行目の要約

詳細な説明をここに書く"

# エディタを開いてコミットメッセージを書く
git commit

# ステージングとコミットを同時に行う（追跡済みファイルのみ）
git commit -am "コミットメッセージ"

# 直前のコミットメッセージを修正する（まだpushしていない場合のみ）
git commit --amend -m "新しいメッセージ"

# ローカルのコミットをリモートリポジトリに送信する
git push

# 初回プッシュ時に上流ブランチを設定する
git push -u origin ブランチ名

# リモートの変更をローカルに取り込む（fetch + merge）
git pull

# リベース形式でリモートの変更を取り込む（履歴がきれいになる）
git pull --rebase
\end{lstlisting}

\section{差分と履歴の確認（→ 第6章、第9章）}

\begin{lstlisting}[language=bash]
# ワーキングディレクトリの変更（未ステージ分）を表示する
git diff

# ステージ済みの変更を表示する
git diff --staged

# 特定のファイルの差分だけ表示する
git diff ファイル名

# ブランチ間の差分を表示する
git diff main..feature-branch

# 変更されたファイル名の一覧だけを表示する
git diff --name-only

# コミット履歴を表示する
git log

# コミット履歴を1行ずつ表示する
git log --oneline

# ブランチの分岐を視覚的に表示する
git log --oneline --graph --all

# 最新N件だけ表示する
git log -5

# 変更内容も含めて表示する
git log -p

# 統計情報（変更行数）を表示する
git log --stat

# 特定ファイルの履歴を表示する
git log -- ファイル名

# 著者で絞り込む
git log --author="名前"

# メッセージで検索する
git log --grep="キーワード"

# ファイルの各行が誰によって変更されたかを表示する
git blame ファイル名

# 特定のコミットの詳細を表示する
git show コミットID
\end{lstlisting}

\section{ブランチ操作（→ 第7章）}

\begin{lstlisting}[language=bash]
# ローカルブランチの一覧を表示する
git branch

# リモートブランチも含めて一覧表示する
git branch -a

# 新しいブランチを作成する（切り替えはしない）
git branch ブランチ名

# ブランチを作成して切り替える
git switch -c ブランチ名

# ブランチを作成して切り替える（従来の書き方）
git checkout -b ブランチ名

# 既存のブランチに切り替える
git switch ブランチ名

# 直前にいたブランチに戻る
git switch -

# マージ済みのブランチを削除する
git branch -d ブランチ名

# 未マージでも強制削除する
git branch -D ブランチ名

# ブランチ名を変更する
git branch -m 新しい名前

# 別のブランチの変更を現在のブランチに統合する
git merge ブランチ名

# マージを中止する
git merge --abort

# 現在のブランチの分岐点を移動させる
git rebase ブランチ名

# リベースを中止する
git rebase --abort

# 特定のコミットだけを現在のブランチに取り込む
git cherry-pick コミットID
\end{lstlisting}

\section{リモート操作（→ 第5章、第6章）}

\begin{lstlisting}[language=bash]
# 登録されているリモートリポジトリの一覧とURLを表示する
git remote -v

# リモートリポジトリを追加する
git remote add origin URL

# リモートリポジトリのURLを変更する
git remote set-url origin 新しいURL

# リモートリポジトリの登録を削除する
git remote remove origin

# リモートの最新情報を取得する（ローカルファイルは変更されない）
git fetch

# 削除されたリモートブランチの参照をクリーンアップする
git fetch --prune

# 特定のブランチをリモートにプッシュする
git push origin ブランチ名

# リモートブランチを削除する
git push origin --delete ブランチ名

# すべてのタグをリモートにプッシュする
git push --tags
\end{lstlisting}

\section{変更の取り消しと復元（→ 第9章）}

\begin{lstlisting}[language=bash]
# ファイルの変更を元に戻す（未ステージの変更を破棄）
git restore ファイル名

# ステージングを解除する（変更内容は保持される）
git restore --staged ファイル名

# コミットを安全に打ち消す（新しいコミットが作られる）
git revert コミットID

# 直前のコミットを取り消す（変更はステージに残る）
git reset --soft HEAD~1

# 直前のコミットとステージングを取り消す（変更はワーキングツリーに残る）
git reset HEAD~1

# 直前のコミットを完全に取り消す（変更が消える。要注意）
git reset --hard HEAD~1

# HEADの移動履歴を表示する（消えたコミットの復旧に使う）
git reflog
\end{lstlisting}

\section{一時退避（→ 第9章）}

\begin{lstlisting}[language=bash]
# 現在の変更を一時退避する
git stash

# メッセージ付きで退避する
git stash save "作業メモ"

# 未追跡ファイルも含めて退避する
git stash -u

# 退避の一覧を表示する
git stash list

# 最新の退避を復元する（退避リストからは削除される）
git stash pop

# 最新の退避を復元する（退避リストに残す）
git stash apply

# 特定の退避を復元する
git stash apply stash@{1}

# 特定の退避を削除する
git stash drop stash@{0}

# すべての退避を削除する
git stash clear
\end{lstlisting}

\section{タグ（→ 第14章）}

\begin{lstlisting}[language=bash]
# タグの一覧を表示する
git tag

# 注釈付きタグを作成する（リリースバージョンに推奨）
git tag -a v1.0.0 -m "バージョン1.0.0"

# 軽量タグを作成する
git tag v1.0.0

# 特定のコミットにタグを付ける
git tag -a v1.0.0 -m "メッセージ" コミットID

# タグを削除する
git tag -d v1.0.0

# タグをリモートにプッシュする
git push origin v1.0.0

# すべてのタグをプッシュする
git push origin --tags

# リモートのタグを削除する
git push origin --delete v1.0.0
\end{lstlisting}

\section{ファイル操作（→ 第6章）}

\begin{lstlisting}[language=bash]
# Gitの管理下でファイル名を変更する
git mv 旧ファイル名 新ファイル名

# Gitの追跡対象からファイルを削除する（ファイル自体も消える）
git rm ファイル名

# Gitの追跡だけを解除する（ファイルは残る）
git rm --cached ファイル名
\end{lstlisting}

\section{GitHub CLI（gh）コマンド（→ 第23章）}

GitHub CLIは、GitHub固有の操作（PR、Issue、Actions、リリースなど）をターミナルから行うためのツールです。\cmd{git} コマンドがローカルのバージョン管理に使われるのに対し、\cmd{gh} コマンドはGitHubプラットフォームとの対話に使われます。

\begin{lstlisting}[language=bash]
# 認証
gh auth login                         # GitHubにログインする
gh auth status                        # 認証状態を確認する

# リポジトリ
gh repo create 名前 --public           # リポジトリを作成する
gh repo clone owner/repo               # リポジトリをクローンする
gh repo list                           # 所有リポジトリの一覧を表示する
gh repo view --web                     # ブラウザでリポジトリを開く

# プルリクエスト
gh pr create --title "タイトル"         # PRを作成する
gh pr list                             # PRの一覧を表示する
gh pr view 番号                         # PRの詳細を表示する
gh pr checkout 番号                     # PRのブランチをローカルに取得する
gh pr merge 番号 --squash               # PRをスカッシュマージする
gh pr diff 番号                         # PRの差分を表示する

# Issue
gh issue create --title "タイトル"      # Issueを作成する
gh issue list                          # Issueの一覧を表示する
gh issue view 番号                      # Issueの詳細を表示する
gh issue close 番号                     # Issueをクローズする

# リリース
gh release create v1.0.0               # リリースを作成する
gh release list                        # リリースの一覧を表示する

# GitHub Actions
gh run list                            # ワークフロー実行の一覧を表示する
gh run view ID                         # 実行の詳細を表示する
gh run rerun ID                        # 失敗した実行を再実行する

# エイリアス
gh alias set prc 'pr create'          # ショートカットを設定する
gh alias list                         # 設定済みエイリアスの一覧を表示する
\end{lstlisting}

\section{逆引きリファレンス --- 「やりたいこと」から探す}

日々の作業でよくある「こうしたい」という場面から、使うべきコマンドを逆引きできる表です。

\begin{table}[h]
\centering
\begin{tabular}{lll}
\hline
やりたいこと & コマンド & 参照 \\
\hline
ファイルの変更を元に戻したい & \cmd{git restore ファイル名} & 第9章 \\
ステージングを取り消したい & \cmd{git restore --staged ファイル名} & 第6章 \\
直前のコミットを取り消したい（変更は残す） & \cmd{git reset --soft HEAD\textasciitilde1} & 第9章 \\
特定のコミットの影響を打ち消したい & \cmd{git revert コミットID} & 第9章 \\
今の作業を一旦脇に置きたい & \cmd{git stash} & 第9章 \\
退避した作業を元に戻したい & \cmd{git stash pop} & 第9章 \\
消してしまったコミットを復元したい & \cmd{git reflog} で履歴を探す & 第9章 \\
別ブランチの特定コミットだけ取り込みたい & \cmd{git cherry-pick コミットID} & 第9章 \\
ファイルの各行を誰が変更したか知りたい & \cmd{git blame ファイル名} & 第9章 \\
リモートリポジトリのURLを変更したい & \cmd{git remote set-url origin URL} & 第5章 \\
\cmd{.gitignore} に追加したのに効かない & \cmd{git rm --cached ファイル名} & 第5章 \\
特定ファイルの変更履歴だけ見たい & \cmd{git log -- ファイル名} & 第6章 \\
間違ったブランチにコミットしてしまった & \cmd{git reset --soft HEAD\textasciitilde1} → \cmd{git stash} → 正しいブランチで \cmd{git stash pop} & 第24章 \\
ブランチ間の差分を確認したい & \cmd{git diff main..feature} & 第6章 \\
リモートの最新状態だけ確認したい & \cmd{git fetch} & 第6章 \\
PRを作りたい & \cmd{gh pr create} & 第10章 \\
Issueを作りたい & \cmd{gh issue create} & 第11章 \\
ワークフローの結果を確認したい & \cmd{gh run list} & 第12章 \\
\hline
\end{tabular}
\end{table}
